% Tensor de estrés en sim-binary-shapes-2d: discos y polígonos regulares.
% Compila con pdflatex, xelatex o lualatex:
%   cd docs/estres && latexmk -pdf estres.tex && latexmk -c
\documentclass[11pt,a4paper]{article}

\usepackage{iftex}
\ifPDFTeX
  \usepackage[T1]{fontenc}
  \usepackage[utf8]{inputenc}
  \usepackage{lmodern}
\else
  \usepackage{fontspec}
\fi
\usepackage[spanish,es-noshorthands,es-tabla]{babel}
\usepackage[margin=2.5cm]{geometry}
\usepackage{amsmath,amssymb}
\usepackage{booktabs}
\usepackage{array}
\usepackage{tikz}
\usetikzlibrary{calc,arrows.meta,positioning}
\usepackage{pgfplots}
\pgfplotsset{compat=1.18}
\usepgfplotslibrary{groupplots}
\usepackage[hidelinks]{hyperref}

\newcommand{\vect}[1]{\mathbf{#1}}
\newcommand{\n}{\vect{n}}
\newcommand{\tv}{\vect{t}}
\newcommand{\nh}{\hat{\vect{n}}}
\newcommand{\tvh}{\hat{\vect{t}}}
\newcommand{\lv}{\vect{l}}
\newcommand{\f}{\vect{f}}
% \code: texto de código que puede cortarse después de cada guion bajo
\newcommand{\codeus}{\textunderscore\allowbreak}
\newcommand{\code}[1]{\texttt{\begingroup\let\_\codeus#1\endgroup}}
\setlength{\emergencystretch}{3em}
\newcommand{\ave}[1]{\langle #1 \rangle}

\title{Tensor de estrés en \code{sim-binary-shapes-2d}:\\
  discos y polígonos regulares}
\author{Manuel Carlevaro}
\date{Rama \code{tensor-en-simulacion}, octubre de 2026}

\begin{document}
\maketitle

\begin{abstract}
  Este documento describe todo el cálculo del tensor de estrés de los
  granos en los atascos del silo: cómo Box2D resuelve los contactos y qué
  se guarda de ellos, la definición del tensor de cada grano y sus
  propiedades para discos y para polígonos regulares, las dos
  implementaciones (en la simulación, con la normal real de cada contacto,
  y en \code{presion\_tensor\_arch.py}, a partir de los archivos), las
  magnitudes derivadas (presión, corte, promedios por región y por
  especie) y la verificación de que ambas coinciden. Cada fórmula remite a
  la función que la implementa.
\end{abstract}

\tableofcontents

% ===========================================================================
\section{Flujo de datos}
\label{sec:flujo}

En cada atasco, antes de romper el arco, el simulador \code{eo-silo}
escribe cuatro archivos en \code{frames\_<dirID>/}
(figura~\ref{fig:flujo}). Las cuatro salidas corresponden al mismo
instante: las posiciones y los impulsos del último paso de integración.

\begin{figure}[htbp]
\centering
\begin{tikzpicture}[>=Stealth, font=\small,
  caja/.style={draw, rounded corners, align=center, inner sep=4pt,
    minimum height=1.0cm},
  arch/.style={draw, align=left, inner sep=3pt, font=\footnotesize\ttfamily,
    fill=gray!8},
  nuevo/.style={arch, fill=blue!8, draw=blue!60!black}]
  \node[caja, minimum width=2.6cm] (sim) at (0,0)
    {\code{eo-silo}\\ atasco detectado};
  \node[arch] (xy) at (5.2,1.95) {atasco\_<jam>.xy};
  \node[arch] (ct) at (5.2,0.65) {atasco\_contacts\_<jam>.dat};
  \node[nuevo] (cf) at (5.2,-0.65) {atasco\_contacts\_full\_<jam>.dat};
  \node[nuevo] (ts) at (5.2,-1.95) {atasco\_tensor\_<jam>.dat};
  \foreach \a in {xy,ct,cf,ts} \draw[->] (sim.east) -- (\a.west);
  \node[caja, minimum width=3.4cm] (py) at (10.6,1.0)
    {\code{presion\_tensor\_arch.py}\\ tensor, $P$, $\tau$, regiones};
  \node[caja, minimum width=3.4cm] (cmp) at (10.6,-1.6)
    {\code{comparar\_tensor\_}\\ \code{simulacion.py}};
  \draw[->] (xy.east) -- (py.west);
  \draw[->] (ct.east) -- (py.west);
  \draw[->, blue!60!black, dashed] (cf.east) -- (py.west);
  \foreach \a in {xy,ct,cf,ts} \draw[->, gray] (\a.east) -- (cmp.west);
  \node[font=\footnotesize, align=left, anchor=north west] at (-1.4,-2.7)
    {En azul: salidas nuevas de la rama \code{tensor-en-simulacion}
     (\code{simulation/src/jam\_output.cpp}).};
\end{tikzpicture}
\caption{Archivos de cada atasco y su uso. \code{saveAtascoFrame} escribe
  la geometría (\code{.xy}), \code{saveAtascoContacts} los contactos en el
  formato de siempre y \code{saveAtascoTensor} el tensor de cada grano y
  los contactos con su normal. El análisis usa la normal de Box2D (flecha
  a trazos) si el archivo \code{atasco\_contacts\_full} está disponible
  (sección~\ref{sec:normal}).}
\label{fig:flujo}
\end{figure}

\begin{center}
\small
\begin{tabular}{>{\raggedright}p{0.33\textwidth}p{0.6\textwidth}}
  \toprule
  Archivo & Contenido \\
  \midrule
  \code{atasco\_<jam>.xy} & Granos activos (\code{isGrain} e \code{isIn}):
    discos como \code{gID 1 x y r tipo}, polígonos como
    \code{gID n x1 y1 ... xn yn tipo}; paredes y tapa. \\
  \code{atasco\_contacts\_<jam>.dat} & Una fila por punto de manifold:
    \code{ID\_A ID\_B PosX PosY NormalImpulse TangentImpulse}, en
    \code{\%.6e}. Sin la normal. \\
  \code{atasco\_contacts\_full\_<jam>.dat} & Las mismas filas, en el mismo
    orden, con la normal de Box2D, las fuerzas, el tipo de manifold y los
    centros. \\
  \code{atasco\_tensor\_<jam>.dat} & Tensor de cada grano calculado en la
    simulación, fuerza y torque netos, contactos y velocidades. \\
  \bottomrule
\end{tabular}
\end{center}

% ===========================================================================
\section{Convenciones}
\label{sec:conv}

\paragraph{Unidades.} Box2D usa unidades MKS. El grano grande de
referencia tiene radio $0.5$, así que su diámetro vale $d = 1$~m, y la
gravedad es $g = 9.8$~m/s$^2$ en la dirección $-y$. El estrés
bidimensional se mide en N/m. \code{presion\_tensor\_arch.py} reporta
además las magnitudes divididas por $g$ (columnas \code{\_reducida}), en
kg/m.

\paragraph{Signo.} Hay dos convenciones, una en cada implementación:
\begin{itemize}
  \item el simulador escribe $s_{ij}$ con \emph{compresión negativa}, como
    la fórmula de Love \eqref{eq:love};
  \item el análisis escribe $\sigma_{ij} = -s_{ij}$, con \emph{compresión
    positiva}.
\end{itemize}
La presión es positiva en compresión en las dos:
$P = -\tfrac12(s_{xx}+s_{yy}) = \tfrac12(\sigma_{xx}+\sigma_{yy})$.

\paragraph{Orden de los índices.} En el simulador,
$s_{ij} \propto \sum f_i\, l_j$: el primer índice va con la fuerza y el
segundo con el brazo. El análisis acumula
$S_{ab} = \sum l_a f_b$, con el brazo primero:
\begin{equation}
  A\, s_{xx} = S_{00}, \qquad A\, s_{xy} = S_{10}, \qquad
  A\, s_{yx} = S_{01}, \qquad A\, s_{yy} = S_{11}.
  \label{eq:mapeo}
\end{equation}
El orden solo importa para la parte antisimétrica.

% ===========================================================================
\section{Geometría de los granos}
\label{sec:geom}

\paragraph{Discos.} Radio $r$ y área $A = \pi r^2$.

\paragraph{Polígonos regulares.} \code{GlobalSetup::load} construye los
$n$ vértices ($3 \le n \le 8$) con el \code{radio} del archivo de
parámetros como radio \emph{circunscripto} $R$:
\begin{equation}
  \vect{v}_k = R\,(\cos 2\pi k/n,\ \sin 2\pi k/n), \qquad k = 0,\dots,n-1.
\end{equation}
Con $\alpha = \pi/n$, la apotema, el lado y el área son
\begin{equation}
  a = R\cos\alpha, \qquad L = 2R\sin\alpha, \qquad
  A = \tfrac{n}{2}R^2\sin 2\alpha = n\,a^2\tan\alpha .
\end{equation}
El centro de masa coincide con el centroide. Para $n \ge 3$, el segundo
momento de área es isótropo: $\int r_i r_j\,dA \propto \delta_{ij}$.

\paragraph{La piel de Box2D.} Un \code{b2PolygonShape} tiene
\code{m\_radius} igual a \code{b2\_polygonRadius}, es decir
$2\times\code{b2\_linearSlop} = r_s = 0.01$~m. El contacto se resuelve
sobre el polígono \emph{redondeado}, inflado en $r_s$
(figura~\ref{fig:geom}a), pero la masa (\code{ComputeMass}) usa solo los
vértices. Los discos no tienen piel.

\paragraph{Área del grano.} Las dos implementaciones usan el área de los
vértices, la misma con la que Box2D calcula la masa:
\begin{itemize}
  \item el simulador la obtiene con \code{ComputeMass(\&md, 1)};
  \item el análisis la calcula con la fórmula del área de Gauss sobre los
    vértices del \code{.xy}.
\end{itemize}
Difieren en $3\times10^{-5}$, por los seis decimales del archivo.

\begin{figure}[htbp]
\centering
\begin{tikzpicture}[>=Stealth]
  % (a) Hexágono con la piel (exagerada)
  \begin{scope}
    \def\R{2}
    \draw[line width=0.36cm, line join=round, gray!35]
      (0:\R) -- (60:\R) -- (120:\R) -- (180:\R) -- (240:\R) -- (300:\R)
      -- cycle;
    \filldraw[thick, fill=white]
      (0:\R) -- (60:\R) -- (120:\R) -- (180:\R) -- (240:\R) -- (300:\R)
      -- cycle;
    \fill (0,0) circle (0.04) node[below left] {$O$};
    \draw (0,0) -- (0:\R) node[midway, below] {$R$};
    \draw (0,0) -- (30:{\R*cos(30)}) node[midway, above left] {$a$};
    \draw (0.55,0) arc[start angle=0, end angle=30, radius=0.55];
    \node at (15:0.8) {$\alpha$};
    \draw ($(30:{\R*cos(30)})+(210:0.2)$) -- ++(120:0.2) -- ++(30:0.2);
    \node at (30:2.25) {$L$};
    \draw[<-] (150:1.82) -- (150:2.55) node[above left] {$r_s$};
    \node at (-2.4,2.1) {(a)};
  \end{scope}
  % (b) Detalle de un contacto cara--cara con piel
  \begin{scope}[shift={(3.3,-0.9)}]
    \begin{scope}
      \clip (0.2,-0.9) rectangle (7.3,2.4);
      \draw[line width=1.8cm, blue!45, draw opacity=0.45, line join=round]
        (-1,0) -- (5,0) -- ++(300:3);
      \draw[line width=1.8cm, red!45, draw opacity=0.45, line join=round]
        (9,1.4) -- (1.5,1.4) -- ++(120:3);
      \fill[white] (-1,0) -- (5,0) -- ++(300:3) -- (-1,-2.6) -- cycle;
      \fill[white] (9,1.4) -- (1.5,1.4) -- ++(120:3) -- (9,4) -- cycle;
      \draw[very thick] (-1,0) -- (5,0) -- ++(300:3);
      \draw[very thick] (9,1.4) -- (1.5,1.4) -- ++(120:3);
      \node at (3.2,-0.5) {A};
      \node at (4.6,1.9) {B};
      \node[font=\scriptsize, blue!50!black] at (0.9,0.35) {piel de A};
      \node[font=\scriptsize, red!50!black] at (6.4,0.95) {piel de B};
      \draw[densely dotted] (1.5,0.7) -- (5,0.7);
      \fill (1.5,0.7) circle (0.07) (5,0.7) circle (0.07);
      \draw[<->] (2.4,0) -- (2.4,0.9) node[midway, left] {\small $r_s$};
      \draw[<->] (3.1,0.5) -- (3.1,0.9) node[midway, right] {\small $\delta$};
      \draw[<->] (4.0,0) -- (4.0,0.7)
        node[midway, right] {\small $r_s - \delta/2$};
      \draw[->, thick, blue] (6.8,-0.8) -- (6.8,0.0) node[right] {$\n$};
    \end{scope}
    \node at (0.5,2.65) {(b)};
    \node at (3.7,-1.3) {$\lambda_A = a + r_s - \delta/2$};
  \end{scope}
\end{tikzpicture}
\caption{(a) Polígono regular ($n = 6$) con radio circunscripto $R$,
  apotema $a$, lado $L$ y $\alpha = \pi/n$. La banda gris es la piel de
  Box2D, de espesor $r_s$, exagerado (en la simulación $r_s = 0.01$);
  los vértices quedan redondeados. (b) Contacto cara--cara entre A
  (abajo) y B (arriba). Las líneas gruesas son las caras y las bandas de
  color las pieles, que se superponen en $\delta$. Box2D ubica los dos
  puntos del manifold (círculos) a mitad de camino entre las pieles, en
  los extremos de la zona en que las caras se enfrentan.}
\label{fig:geom}
\end{figure}

% ===========================================================================
\section{Fuerzas de contacto en Box2D}
\label{sec:box2d}

Box2D 2.4.2 resuelve los contactos con impulsos. En cada paso $\Delta t$,
el solver de velocidades acumula en cada punto $i$ del manifold un impulso
normal $\lambda_{n,i} \ge 0$ (\code{normalImpulse}) y uno tangencial
$\lambda_{t,i}$ (\code{tangentImpulse}), con $|\lambda_t| \le \mu
\lambda_n$. Las fuerzas medias durante el paso son
\begin{equation}
  F_n = \lambda_n/\Delta t, \qquad F_t = \lambda_t/\Delta t .
\end{equation}

\paragraph{Normal y tangente.} La normal $\n$ del manifold apunta del
cuerpo A al cuerpo B (A y B son las \emph{fixtures} del contacto). La
tangente es la que usa \code{b2ContactSolver}:
\begin{equation}
  \tv = \code{b2Cross}(\n, 1) = (n_y, -n_x).
\end{equation}
La fuerza sobre cada cuerpo es
\begin{equation}
  \f_B = F_n\,\n + F_t\,\tv, \qquad \f_A = -\f_B .
  \label{eq:fuerza}
\end{equation}
El análisis escribe la fuerza sobre A con $\tv' = (-n_y, n_x) = -\tv$:
$\f_A = (-F_n\,\n + F_t\,\tv')$ (\code{fuerzas\_contacto}). Es la misma
fuerza. Se validó con el equilibrio estático de cada grano: el residuo es
$\sim 10^{-6}$, mientras que con la tangente invertida es del orden de
$mg$ (sección~\ref{sec:verif}).

\paragraph{Tipos de manifold.} El campo \code{type} del manifold (columna
\code{mtype} de \code{atasco\_contacts\_full}) distingue tres casos:
\begin{itemize}
  \item \code{e\_circles} (0): disco con disco. Un punto, sobre la línea
    de los centros.
  \item \code{e\_faceA} (1) y \code{e\_faceB} (2): hay un polígono o una
    pared. La normal es la de una \emph{cara de referencia} de A o de B.
  \item Contra un polígono, el manifold tiene uno o dos puntos: dos en un
    contacto cara--cara (figura~\ref{fig:geom}b) y uno en un contacto de
    vértice o contra un disco.
\end{itemize}
Los contactos con las paredes (cadena del silo) tienen un manifold por
arista. Por eso un mismo par grano--pared puede aparecer en varias filas.

\paragraph{Ubicación del punto.} Box2D ubica cada punto a mitad de camino
entre las superficies de los dos cuerpos, piel incluida. Para un disco,
el punto queda sobre la recta normal que pasa por su centro. Para la cara
de referencia de un polígono, queda a una distancia $r_s - \delta/2$ de
la cara (figura~\ref{fig:geom}b).

\paragraph{Qué impulsos se guardan.} Al escribir el atasco, los impulsos
del manifold son los del último \code{Step}. Si los cuerpos de un contacto
están dormidos, Box2D no lo vuelve a resolver y conserva los impulsos del
último paso en que lo resolvió. En un atasco en reposo esos impulsos
siguen equilibrando el peso: el residuo de equilibrio no cambia
(sección~\ref{sec:verif}).

% ===========================================================================
\section{Tensor de estrés de cada grano}
\label{sec:tensor}

\subsection{Definición}

Para cada grano $p$ se suma sobre todos sus puntos de contacto $c$, con
otros granos y con las paredes:
\begin{equation}
  s^{(p)}_{ij} = \frac{1}{A_p}\sum_c f^{(c)}_i\, l^{(c)}_j, \qquad
  s^{n(p)}_{ij} = \frac{1}{A_p}\sum_c f^{n(c)}_i\, l^{(c)}_j, \qquad
  s^t = s - s^n,
  \label{eq:love}
\end{equation}
donde:
\begin{itemize}
  \item $\f^{(c)}$ es la fuerza sobre el grano según \eqref{eq:fuerza};
  \item $\f^{n(c)} = \pm F_n\n$ es solo su componente normal;
  \item $\lv^{(c)}$ es el vector del centro de masa al punto de contacto;
  \item $A_p$ es el área del grano.
\end{itemize}

\paragraph{Relación con el estrés de Cauchy.} Para un cuerpo rígido, el
teorema de la divergencia da
\begin{equation}
  \int_{A_p}\sigma^{\text{C}}_{ij}\,dA = \sum_c f_i l_j
    + \int_{A_p} \rho\, g_i\, r_j\,dA - \int_{A_p}\rho\, a_i\, r_j\,dA .
\end{equation}
La gravedad es uniforme y $\int \vect{r}\,dA = 0$ respecto del centro de
masa, así que no aporta. En un atasco en reposo, la aceleración es nula.
Por lo tanto, $s^{(p)}$ es el estrés medio de Cauchy del grano: la
fórmula de Love vale para cualquier forma.

\subsection{Brazo de palanca y marco del grano}

Para cada grano y cada punto de contacto se definen:
\begin{itemize}
  \item la normal saliente $\nh$: $\nh = \n$ para A y $\nh = -\n$ para B;
  \item la tangente $\tvh = (\hat n_y, -\hat n_x)$.
\end{itemize}
Con \eqref{eq:fuerza}, la fuerza sobre el grano es, en los dos casos,
\begin{equation}
  \f = -F_n\,\nh - F_t\,\tvh ,
\end{equation}
y el brazo se descompone como
\begin{equation}
  \lv = \lambda\,\nh + \eta\,\tvh, \qquad \lambda = \lv\cdot\nh, \quad
  \eta = \lv\cdot\tvh .
\end{equation}
$\lambda$ es la distancia del centro a la recta de contacto y $\eta$ la
\emph{excentricidad} del punto (figura~\ref{fig:contactos}).

\begin{itemize}
  \item Disco: $\lambda = r - \delta/2$ y $\eta = 0$.
  \item Polígono, contacto cara--cara: para el dueño de la cara de
    referencia, $\lambda \simeq a + r_s - \delta/2$ y $|\eta| \le L/2$.
  \item Polígono, contacto de vértice: para el grano que apoya el
    vértice, $\lambda < R$ y $|\eta|$ puede ser del orden de $R$.
\end{itemize}

\begin{figure}[htbp]
\centering
\begin{tikzpicture}[scale=1.1, >=Stealth]
  % (a) Disco
  \begin{scope}
    \draw[thick] (0,0) circle (1);
    \draw[thick] (1.9,0) circle (1);
    \fill (0,0) circle (0.04);
    \node at (-0.45,-0.45) {$p$};
    \node at (2.35,-0.45) {$q$};
    \draw[->, gray, thick] (0,0) -- (0.92,0) node[midway, above] {$\lv$};
    \fill (0.95,0) circle (0.05);
    \draw[->, very thick, blue] (0.95,0) -- (1.6,0) node[above] {$\nh$};
    \draw[->, very thick, red] (0.95,0) -- (0.95,-0.65)
      node[below] {$\tvh$};
    \node at (-1.1,1.35) {(a)};
    \node at (0.95,-1.65) {$\lambda = r - \delta/2, \ \eta = 0$};
  \end{scope}
  % (b) Cara--cara
  \begin{scope}[shift={(4.7,0)}]
    \draw[thick] (-1,-1) rectangle (1,1);
    \draw[thick] (1,-0.55) rectangle (3,1.45);
    \node at (-0.6,-0.7) {$p$};
    \node at (2.5,0.95) {$q$};
    \fill (0,0) circle (0.04);
    \draw[->, gray, thick] (0,0) -- (0.96,0.96)
      node[pos=0.6, above left] {$\lv_1$};
    \draw[->, gray, thick] (0,0) -- (0.95,-0.52)
      node[pos=0.5, below left] {$\lv_2$};
    \draw[dashed] (0,0) -- (1,0);
    \node[above] at (0.8,0) {$\lambda$};
    \fill (1,1) circle (0.06) node[above right] {1};
    \fill (1,-0.55) circle (0.06) node[below right] {2};
    \draw[densely dotted] (1,0) -- (1.5,0);
    \draw[<->] (1.35,0) -- (1.35,1) node[midway, right] {$\eta_1$};
    \draw[<->] (1.35,0) -- (1.35,-0.55) node[midway, right] {$\eta_2$};
    \draw[->, very thick, blue] (-0.85,0.75) -- (-0.4,0.75)
      node[right] {$\nh$};
    \draw[->, very thick, red] (-0.85,0.75) -- (-0.85,0.3)
      node[right] {$\tvh$};
    \node at (-1.1,1.35) {(b)};
    \node at (1,-1.65) {dos puntos, un contacto};
  \end{scope}
  % (c) Vértice--cara
  \begin{scope}[shift={(9.7,0)}]
    \def\Rc{1.4142}
    \draw[thick] (25:\Rc) -- (115:\Rc) -- (205:\Rc) -- (295:\Rc) -- cycle;
    \draw[thick] (1.2817,-1) rectangle (2.9,1.25);
    \node at (0.3,-0.6) {$p$};
    \node at (2.45,-0.65) {$q$};
    \fill (0,0) circle (0.04);
    \draw[->, gray, thick] (0,0) -- (25:1.37)
      node[pos=0.5, above] {$\lv$};
    \draw[dashed] (0,0) -- (1.2817,0) node[midway, below] {$\lambda$};
    \fill (25:\Rc) circle (0.06);
    \draw[densely dotted] (1.2817,0) -- (1.8,0);
    \draw[<->] (1.65,0) -- (1.65,0.5977) node[midway, right] {$\eta$};
    \draw[->, very thick, blue] (-0.6,-0.15) -- (-0.15,-0.15)
      node[right] {$\nh$};
    \draw[->, very thick, red] (-0.6,-0.15) -- (-0.6,-0.6)
      node[right] {$\tvh$};
    \node at (-1.3,1.35) {(c)};
    \node at (0.8,-1.65) {$\lambda < R$, \ $|\eta| \sim R$};
  \end{scope}
\end{tikzpicture}
\caption{Marco del grano $p$ y descomposición del brazo
  $\lv = \lambda\nh + \eta\tvh$. $\nh$ es la normal saliente de $p$
  (hacia $q$) y $\tvh = (\hat n_y, -\hat n_x)$, hacia abajo. (a) Disco:
  el punto está sobre la línea de los centros y $\lv\parallel\nh$.
  (b) Contacto cara--cara: el manifold tiene dos puntos con el mismo
  $\lambda$ y excentricidades distintas ($\eta_1 < 0$, $\eta_2 > 0$,
  medidas sobre $\tvh$). Es \emph{un} contacto con dos puntos. (c)
  Contacto vértice--cara: $\nh$ es la normal de la cara de $q$ y $\lv$
  apunta al vértice de $p$. Se omite la piel.}
\label{fig:contactos}
\end{figure}

\subsection{Formas explícitas y propiedades}
\label{sec:prop}

Sustituyendo en \eqref{eq:love},
\begin{align}
  A_p\, s^{n(p)} &= -\sum_c F_n\left(\lambda\,\nh\otimes\nh
                     + \eta\,\nh\otimes\tvh\right), \label{eq:sn}\\
  A_p\, s^{t(p)} &= -\sum_c F_t\left(\lambda\,\tvh\otimes\nh
                     + \eta\,\tvh\otimes\tvh\right). \label{eq:st}
\end{align}
Con $\nh\times\tvh = -1$ y
$(\vect{a}\otimes\vect{b})_{yx} - (\vect{a}\otimes\vect{b})_{xy} =
-\vect{a}\times\vect{b}$ se obtienen la presión y la parte antisimétrica:
\begin{align}
  P^{(p)} &= -\tfrac12\operatorname{tr}s^{(p)}
    = \frac{1}{2A_p}\sum_c\left(F_n\lambda + F_t\eta\right),
    \label{eq:presion}\\
  s_{yx} - s_{xy} &= \frac{1}{A_p}\sum_c\left(F_t\lambda - F_n\eta\right)
    = \frac{\tau^t_\text{c} + \tau^n_\text{c}}{A_p}
    = \frac{\tau_\text{c}}{A_p},
    \label{eq:torque}
\end{align}
donde $\tau_\text{c} = \sum_c \lv\times\f$ es el torque de contacto
respecto del centro de masa. La tabla resume las propiedades para cada
forma.

\begin{center}
\small
\begin{tabular}{>{\raggedright}p{0.36\textwidth}>{\raggedright}p{0.24\textwidth}p{0.30\textwidth}}
  \toprule
  Propiedad & Discos ($\eta = 0$) & Polígonos \\
  \midrule
  $s^n$ simétrico & sí & no: $s^n_{yx}-s^n_{xy} = -\sum F_n\eta/A_p$ \\
  $\operatorname{tr}s^t = 0$ (la fricción no aporta a $P$)
    & sí & no: $\operatorname{tr}s^t = -\sum F_t\eta/A_p$ \\
  $A_p(s_{yx}-s_{xy}) = \tau_\text{c}$ & sí & sí \\
  En reposo ($\tau_\text{c} = 0$): $s$ simétrico & sí & sí \\
  Presión por cara: $P^n \simeq \sum F_n/P_\text{er}$
    & $P_\text{er} = 2\pi r$ & $P_\text{er} = nL$ (contactos cara--cara) \\
  \bottomrule
\end{tabular}
\end{center}

En un atasco en reposo, el tensor total de cada grano es simétrico, como
se afirma en el artículo. Con polígonos, en cambio, sus partes normal y
tangencial no lo son por separado: los dos torques $\tau^n_\text{c}$ y
$\tau^t_\text{c}$ pueden ser grandes y cancelarse
(figura~\ref{fig:torque}).

\begin{figure}[htbp]
\centering
\begin{tikzpicture}[scale=1.25, >=Stealth]
  \begin{scope}
    \draw[thick] (-1.4,-1.4) rectangle (1.4,1.4);
    \fill (0,0) circle (0.045) node[above left] {$O$};
    \draw[->, gray, thick] (0,0) -- (1.36,0.875)
      node[pos=0.6, above left] {$\lv$};
    \draw[dashed] (0,0) -- (1.4,0) node[pos=0.4, below] {$\lambda$};
    \draw[densely dotted] (1.4,0) -- (1.95,0);
    \draw[<->] (1.75,0) -- (1.75,0.9) node[midway, right] {$\eta$};
    \draw[->, very thick, blue!70!black] (1.4,0.9) -- (0.35,0.9)
      node[pos=0.6, above] {$-F_n\nh$};
    \draw[->, very thick, red!70!black] (1.4,0.9) -- (1.4,1.95)
      node[above] {$-F_t\tvh$};
    \fill (1.4,0.9) circle (0.07);
    \draw[->, very thick, blue] (-1.1,-0.6) -- (-0.6,-0.6)
      node[right] {$\nh$};
    \draw[->, very thick, red] (-1.1,-0.6) -- (-1.1,-1.1)
      node[right] {$\tvh$};
    \node[align=left, anchor=west] at (2.5,0.2)
      {$\tau^n = -F_n\,\eta$\\[3pt]
       $\tau^t = F_t\,\lambda$\\[3pt]
       $2A_p\,P = F_n\lambda + F_t\eta$};
    \node at (-1.6,1.9) {(a)};
  \end{scope}
  \begin{scope}[shift={(8.5,0)}]
    \draw[thick] (-1.4,-1.4) rectangle (1.4,1.4);
    \fill (0,0) circle (0.045);
    \draw[->, very thick, blue!70!black] (2.4,0.8) -- (1.4,0.8)
      node[pos=0.4, above] {$F_n$};
    \draw[->, very thick, blue!70!black] (-2.4,-0.8) -- (-1.4,-0.8)
      node[pos=0.4, below] {$F_n$};
    \fill (1.4,0.8) circle (0.07) (-1.4,-0.8) circle (0.07);
    \draw[->, thick] (-30:0.5) arc[start angle=-30, end angle=210,
      radius=0.5];
    \node at (0,0.8) {$\tau^n > 0$};
    \node at (-2.3,1.9) {(b)};
  \end{scope}
\end{tikzpicture}
\caption{(a) Fuerza sobre el grano en un punto de su cara derecha ($\nh$
  hacia la derecha, $\tvh$ hacia abajo). Como el punto está por encima
  del centro, $\eta < 0$: la fuerza normal produce un torque antihorario
  $\tau^n = -F_n\eta > 0$, y la tangencial aporta a la presión con
  $F_t\eta$. En un disco, $\eta = 0$ y ninguno de los dos términos
  existe. (b) Dos fuerzas puramente normales con resultante nula: el
  grano está en equilibrio de fuerzas pero recibe un torque, y
  $s^n_{yx}-s^n_{xy} \ne 0$. En un disco, toda fuerza normal pasa por el
  centro.}
\label{fig:torque}
\end{figure}

% ===========================================================================
\section{La normal del contacto}
\label{sec:normal}

El tensor depende de la normal de cada contacto, que separa la fuerza en
sus partes normal y tangencial \eqref{eq:fuerza}. Hay dos maneras de
obtenerla.

\subsection{Normal de Box2D}

La normal del manifold es la que usó el solver
(\code{b2WorldManifold::}\allowbreak\code{normal}).
\begin{itemize}
  \item En la simulación, \code{saveAtascoTensor} la usa directamente y
    la escribe en \code{atasco\_contacts\_full}.
  \item En el análisis, \code{leer\_normales\_box2d} la lee cuando ese
    archivo existe y coincide fila por fila con \code{atasco\_contacts}.
\end{itemize}
La columna \code{normales} del resumen por atasco indica cuál se usó:
\code{box2d} o \code{inferidas (...)}.

\subsection{Normal inferida de la geometría}

Con los datos anteriores a la rama \code{tensor-en-simulacion} no hay
normal guardada. \code{normales\_contacto} la infiere, orientada de A
hacia B:
\begin{itemize}
  \item \emph{Si uno de los cuerpos es un disco}, usa la dirección radial
    del disco, del centro al punto de contacto. Es exacta.
  \item \emph{Polígono contra polígono o contra pared}: para cada lado
    busca la cara más cercana al punto de contacto (\code{normales\_poligonos},
    \code{normales\_paredes}) y su \emph{interioridad}, la distancia de
    la proyección del punto al vértice más cercano de esa cara. Usa la
    cara con mayor interioridad, que debería ser la cara de referencia de
    Box2D.
\end{itemize}

\paragraph{Dónde falla.} En un contacto \emph{esquina contra esquina}
entre dos polígonos, el punto de Box2D queda a $\sim r_s$ de un vértice
de \emph{cada} polígono (figura~\ref{fig:esquina}). Las cuatro caras
candidatas tienen interioridad $\approx 0$ y la elección es arbitraria.
Si se elige la cara de al lado, la normal sale girada $90^\circ$ (en un
cuadrado). Esa fuerza se descompone mal y el tensor de los dos granos
queda mal. La posición del punto no alcanza para resolver la ambigüedad:
hace falta la normal que usó el solver.

\begin{figure}[htbp]
\centering
\begin{tikzpicture}[>=Stealth, scale=1.3]
  % Núcleos y pieles de dos cuadrados que se tocan por las esquinas
  \begin{scope}
    \clip (-2.6,-2.3) rectangle (3.4,2.6);
    \draw[line width=1.17cm, blue!45, draw opacity=0.45, line join=round]
      (-3,0) -- (0,0) -- (0,-3);
    \draw[line width=1.17cm, red!45, draw opacity=0.45, line join=round]
      (4,0.55) -- (0.62,0.55) -- (0.62,4);
    \fill[white] (-3,0) -- (0,0) -- (0,-3) -- (-3,-3) -- cycle;
    \fill[white] (4,0.55) -- (0.62,0.55) -- (0.62,4) -- (4,4) -- cycle;
    \draw[very thick] (-3,0) -- (0,0) -- (0,-3);
    \draw[very thick] (4,0.55) -- (0.62,0.55) -- (0.62,4);
  \end{scope}
  \node at (-1.4,-1.3) {A};
  \node at (2.0,1.8) {B};
  \fill (0.31,0.275) circle (0.06);
  \draw[gray] (0.31,0.275) -- (-0.35,1.05)
    node[above left, font=\small, inner sep=1pt] {punto de contacto};
  \draw[->, ultra thick, red!70!black] (0.31,0.275) -- ++(1.2,0)
    node[below right, font=\small, align=left] {Box2D (cara de A)};
  \draw[->, ultra thick, blue!70!black, dashed] (0.31,0.275) -- ++(0,1.2)
    node[above, font=\small] {inferida};
  \node[font=\scriptsize, blue!50!black] at (-1.5,0.22) {piel de A};
  \node[font=\scriptsize, red!50!black] at (2.9,0.33) {piel de B};
\end{tikzpicture}
\caption{Contacto esquina contra esquina entre dos cuadrados, con la piel
  exagerada. El punto de Box2D (círculo) queda junto a un vértice de cada
  polígono, así que las cuatro caras son igualmente candidatas. En este
  ejemplo, Box2D eligió la cara derecha de A (flecha llena) y la
  inferencia, la cara superior (flecha a trazos): la normal queda girada
  $90^\circ$.}
\label{fig:esquina}
\end{figure}

% ===========================================================================
\section{Implementación en la simulación}
\label{sec:sim}

\code{saveAtascoTensor} (\code{simulation/src/jam\_output.cpp}) recorre
la lista de contactos de Box2D que se tocan. Para cada punto $i$ de cada
manifold:
\begin{enumerate}
  \item calcula $F_n$, $F_t$, $\n$ y $\tv$ y la fuerza
    \eqref{eq:fuerza} sobre A y sobre B;
  \item a cada cuerpo que sea un grano activo le suma $\f\otimes\lv$ (y la
    parte normal), la fuerza y el torque $\lv\times\f$, con
    $\lv = $ punto $-$ \code{GetWorldCenter()};
  \item escribe la fila del punto en \code{atasco\_contacts\_full}.
\end{enumerate}
Al final divide por el área (\code{ComputeMass}) y escribe una fila por
grano activo en \code{atasco\_tensor}:
\begin{center}
  \small\code{gID tipo nLados x y A m sxx sxy syx syy snxx snxy snyx snyy
  Fx Fy tau z z\_w z\_f vx vy w}
\end{center}

\paragraph{Contactos.} \code{z} cuenta \emph{contactos} (manifolds con
$\sum_i F_{n,i} > 0$), no puntos: un contacto cara--cara cuenta una vez.
\code{z\_w} es la parte de ellos con paredes y \code{z\_f} la parte con
los dos puntos activos (cara--cara). La fracción $z_f/z$ indica cuánto de
la red de fuerzas se transmite por caras, que es lo que controla $\eta$.

\paragraph{Velocidades.} \code{vx}, \code{vy} y \code{w} permiten
comprobar que el atasco está en reposo.

% ===========================================================================
\section{Implementación en el análisis}
\label{sec:py}

\subsection{Tensor, presión y corte de cada grano}

\code{calcular\_tensor\_y\_presion\_vectorizado} hace los siguientes
pasos:
\begin{enumerate}
  \item une cada contacto con los dos granos del \code{.xy}: centroide,
    forma y vértices;
  \item obtiene la normal (sección~\ref{sec:normal}) y la fuerza sobre A
    (\code{fuerzas\_contacto});
  \item acumula $S_{ab} = \sum_c l_a f_b$ en A con $\f_A$ y en B con
    $-\f_A$. Los contactos con paredes (\code{id} $<0$) aportan solo al
    grano;
  \item calcula, con compresión positiva,
    \begin{equation}
      P = -\frac{S_{00}+S_{11}}{2A}, \quad
      \sigma_{xx} = -\frac{S_{00}}{A}, \quad
      \sigma_{yy} = -\frac{S_{11}}{A}, \quad
      \sigma_{xy} = -\frac{S_{01}+S_{10}}{2A},
    \end{equation}
    \begin{equation}
      \tau = \sqrt{\left(\frac{\sigma_{xx}-\sigma_{yy}}{2}\right)^2
                   + \sigma_{xy}^2}.
    \end{equation}
\end{enumerate}
$\sigma_{xy}$ es la parte simétrica del tensor. La antisimétrica, nula en
reposo, no se usa. $\tau$ es el radio del círculo de Mohr: el corte
máximo. Las columnas \code{\_reducida} son las mismas magnitudes
divididas por $g$.

\subsection{Especies, arco y regiones}

\paragraph{Especies.} \code{especie\_por\_tipo} (\code{analysis/especies.py})
clasifica cada grano por su columna \code{tipo} del \code{.xy}, no por un
umbral de radio:
\begin{itemize}
  \item la especie de mayor área mediana es \code{B} (grande);
  \item la otra es \code{S} (chica);
  \item con $\chi = 0$, con una sola especie, o con áreas iguales dentro
    del 1\,\%, todos los granos son \code{B}.
\end{itemize}

\paragraph{Arco.} \code{detectar\_arco\_primario} arma el grafo de
contactos grano--grano: una arista por par, con costo $10^{y_\text{max}}$,
donde $y_\text{max}$ es la altura del más alto de los dos centros. Busca
el camino de costo mínimo (Dijkstra) entre un grano apoyado a la
izquierda del orificio y uno apoyado a la derecha. Ese costo favorece
fuertemente los caminos bajos. Después recorta de las puntas los granos
que son solo cimiento: apoyados en el piso fuera del orificio, o con dos o
más apoyos inferiores ajenos al arco.

\paragraph{Regiones.} Hay tres (figura~\ref{fig:regiones}):
\begin{itemize}
  \item \emph{arco}: los granos del camino;
  \item \emph{ROI}: los demás granos con centro en la franja horizontal
    $y^\text{arco}_\text{min} \le y \le y^\text{arco}_\text{max} +
    \code{ROI\_EXTRA}$, de pared a pared, con $\code{ROI\_EXTRA} = 0$ por
    defecto;
  \item \emph{bulk}: el resto.
\end{itemize}

\begin{figure}[htbp]
\centering
\begin{tikzpicture}[scale=0.75]
  \fill[orange!15] (-5,0.75) rectangle (5,2.0);
  \draw[very thick] (-5,4.2) -- (-5,0) -- (-1.1,0);
  \draw[very thick] (1.1,0) -- (5,0) -- (5,4.2);
  \foreach \x/\y in {-3.5/0.5,-2.5/0.5,2.6/0.5,3.6/0.5,-4.5/0.5,4.5/0.5,
                     -4.4/1.5,-3.4/1.6,3.3/1.55,4.4/1.5,-4.4/2.6,-3.3/2.7,
                     -2.2/2.7,-1.1/3.0,0/3.1,1.1/2.95,2.2/2.7,3.3/2.7,
                     4.4/2.6,-2.6/3.7,-1.5/3.9,0.5/4.0,1.6/3.9,2.8/3.7}
    \draw[gray] (\x,\y) circle (0.48);
  \foreach \x/\y in {-1.65/0.75,-0.75/1.55,0.3/2.0,1.35/1.6,2.0/0.8}
    \filldraw[fill=red!30, draw=red!60!black, thick] (\x,\y) circle (0.48);
  \foreach \x/\y in {-2.6/1.5,2.4/1.75}
    \filldraw[fill=blue!20, draw=blue!60!black] (\x,\y) circle (0.48);
  \draw[dashed] (-5.4,0.75) -- (5.4,0.75) node[right, font=\small]
    {$y^\text{arco}_\text{min}$};
  \draw[dashed] (-5.4,2.0) -- (5.4,2.0) node[right, font=\small]
    {$y^\text{arco}_\text{max}$};
  \node[font=\small, red!60!black] at (0,-0.6) {arco};
  \node[font=\small, blue!60!black, align=center] at (-6.6,1.4)
    {ROI\\[-1pt]\scriptsize(franja,\\[-2pt]\scriptsize sin el arco)};
  \node[font=\small] at (-6.6,3.3) {bulk};
  \draw[<->] (-1.1,-0.25) -- (1.1,-0.25);
  \node[font=\scriptsize] at (0,-1.1) {orificio $2r$};
\end{tikzpicture}
\caption{Regiones de un atasco. En rojo, el arco: camino de contactos de
  costo mínimo entre los dos bordes del orificio, sin los cimientos. En
  azul, granos de la ROI: los que no son del arco y tienen el centro en la
  franja (sombreada) entre los centros más bajo y más alto del arco, de
  pared a pared. El resto es el bulk.}
\label{fig:regiones}
\end{figure}

\subsection{Promedios}

Para cada atasco, región $R$ y especie $e$ se guardan los numeradores y
denominadores de los promedios pesados por área:
\begin{equation}
  \code{SumP\_}R\code{\_}e = \sum_{p\in R,e} P_p A_p, \qquad
  \code{SumTau\_}R\code{\_}e = \sum_{p\in R,e} \tau_p A_p, \qquad
  \code{Area\_}R\code{\_}e = \sum_{p\in R,e} A_p .
\end{equation}
\code{resumen\_por\_configuracion} los suma sobre todos los atascos de una
configuración (forma, $\chi$, razón de tamaños) y divide:
\begin{equation}
  \ave{P}_{R,e} = \frac{\sum_\text{atascos}\code{SumP\_}R\code{\_}e}
                       {\sum_\text{atascos}\code{Area\_}R\code{\_}e},
\end{equation}
y lo mismo para $\ave{\tau}_{R,e}$. Es el estrés medio sobre los granos
de esa especie y región: cada atasco pesa según el área de granos que
aporta, no por igual.

% ===========================================================================
\section{Verificación}
\label{sec:verif}

\paragraph{Prueba.} Simulación local de un silo de semiancho 5 y alto
40, con 250 cuadrados ($R = 0.5$) y 120 discos ($r = 0.158$), orificio
$r = 1.1$ y \code{tEsperaAtasco: 5.0}, como en producción. Se analizaron
tres atascos, con 1110 pares grano--atasco. El script es
\code{analysis/utils/comparar\_tensor\_simulacion.py}.

\paragraph{Controles internos del tensor del simulador.} Se cumplen
cuatro identidades:
\begin{itemize}
  \item \emph{Equilibrio de fuerzas}: $|\vect{F}+m\vect{g}|/mg$ tiene
    una mediana de $\sim 10^{-6}$ y un máximo de $1.2\times10^{-4}$.
  \item \emph{Equilibrio de torque}: $|\tau_\text{c}|/(mgR)$ tiene una
    mediana de $\sim 4\times10^{-7}$ y un máximo de $7\times10^{-5}$. El
    tensor total es simétrico.
  \item \emph{Identidad \eqref{eq:torque}}: $A(s_{yx}-s_{xy}) =
    \tau_\text{c}$ a $4\times10^{-8}$.
  \item \emph{Virial global}: sumando sobre los granos, los contactos
    internos se cancelan y queda una identidad exacta que controla
    signos, brazos y puntos para cualquier forma:
    \begin{equation}
      \sum_p A_p s^{(p)} + \sum_p \vect{F}_p\otimes\vect{x}_p
      = \sum_{c\in\text{ext}}\f_c\otimes\vect{x}_c .
    \end{equation}
    Se cumple a $10^{-10}$.
\end{itemize}

\paragraph{Comparación con el análisis.} La tabla compara, para cada
grano, la presión del análisis con la de la simulación. El error es
relativo a $|P| + q$ del simulador, con $q$ el corte máximo.

\begin{center}
\small
\begin{tabular}{lccc}
  \toprule
  Normal usada por el análisis & mediana & granos $>1$\,\% & granos $>10$\,\% \\
  \midrule
  inferida (datos viejos) & $3\times10^{-6}$ & 34 (3.1\,\%) & 21 (1.9\,\%) \\
  Box2D (\code{atasco\_contacts\_full}) & $1\times10^{-6}$ & 0 & 0 \\
  \bottomrule
\end{tabular}
\end{center}

Con la normal inferida:
\begin{itemize}
  \item los discos y las paredes salen exactos;
  \item entre el 1 y el 2.5\,\% de los puntos de contacto entre
    polígonos tiene la normal girada $90^\circ$; son todos contactos esquina contra
    esquina (figura~\ref{fig:esquina});
  \item las medias por región ($\ave{P}$ en el arco y en la ROI)
    cambiaron $\le 0.03$\,\% en estos tres atascos, pero el valor de
    cada grano afectado, por ejemplo en los mapas de presión, está mal.
\end{itemize}
Con la normal de Box2D, el análisis coincide con la simulación en todos
los granos: el error máximo es $6\times10^{-5}$, la precisión de los
archivos (figura~\ref{fig:comparacion}).

\begin{figure}[htbp]
\centering
\begin{tikzpicture}
  \begin{groupplot}[group style={group size=2 by 1, horizontal sep=1.6cm},
      width=6.6cm, height=6.6cm, xmode=log, ymode=log,
      xmin=0.5, xmax=2000, ymin=0.5, ymax=2000,
      xlabel={$P$ simulación (N/m)},
      tick label style={font=\footnotesize}, label style={font=\small},
      title style={font=\small}]
    \nextgroupplot[title={Normal inferida}, ylabel={$P$ análisis (N/m)}]
      \addplot[gray, domain=0.5:2000, samples=2] {x};
      \addplot[only marks, mark=*, mark size=0.8pt, blue!60!black,
        restrict expr to domain={\thisrow{err}}{0:0.01}]
        table[x=P_sim, y=P_py] {datos/p_inferidas.dat};
      \addplot[only marks, mark=*, mark size=1.6pt, red!80!black,
        restrict expr to domain={\thisrow{err}}{0.01:1e9}]
        table[x=P_sim, y=P_py] {datos/p_inferidas.dat};
    \nextgroupplot[title={Normal de Box2D}]
      \addplot[gray, domain=0.5:2000, samples=2] {x};
      \addplot[only marks, mark=*, mark size=0.8pt, blue!60!black]
        table[x=P_sim, y=P_py] {datos/p_auto.dat};
  \end{groupplot}
\end{tikzpicture}
\caption{Presión de cada grano en los tres atascos de prueba, calculada
  por el análisis frente a la calculada en la simulación. En rojo, los
  granos con error mayor que 1\,\%: todos tienen algún contacto de
  esquina entre polígonos con la normal mal inferida. Con la normal de Box2D,
  todos los puntos caen sobre la diagonal.}
\label{fig:comparacion}
\end{figure}

\paragraph{Atascos que no están en reposo.} Con
\code{tEsperaAtasco: 1.0}, se guardó un atasco 1~s después de un
desatasco que no reanudó el flujo. El empaquetamiento todavía se estaba
reacomodando: la mediana de $|\vect{F}+m\vect{g}|/mg$ era $0.66$. Con los
5~s de producción no ocurrió en la prueba, pero es un caso posible.
Conviene descartar esos atascos con las columnas de velocidad o con el
residuo de equilibrio.

\paragraph{Pruebas sintéticas.} \code{analysis/utils/verificar\_tensor.py
--sintetico} verifica el signo de $\sigma_{xy}$, la parte tangencial, la
tercera ley y la invariancia ante rotaciones con casos de resultado
conocido. Siguen dando correcto con los cambios de la rama.

% ===========================================================================
\section{Resumen y recomendaciones}

\begin{itemize}
  \item La fórmula de Love \eqref{eq:love} vale para discos y polígonos.
    En reposo, el tensor total de cada grano es simétrico. Con polígonos,
    $s^n$ y $s^t$ no lo son por separado, y la fricción aporta a la
    presión.
  \item Las dos implementaciones coinciden a la precisión de los
    archivos cuando el análisis usa la normal de Box2D.
  \item Con los datos existentes, sin \code{atasco\_contacts\_full}, la
    normal inferida falla en los contactos esquina contra esquina entre
    polígonos. Afecta los valores de un 2--3\,\% de los granos y casi no
    afecta las medias por región. Para los mapas por grano conviene
    volver a correr con el simulador de esta rama.
  \item Conviene filtrar los atascos que no están en reposo.
  \item Las presiones de especies distintas se comparan a igual
    normalización por el área del grano: a igual fuerza, el grano chico
    tiene más presión. Para separar fuerza y tamaño conviene reportar
    también $A_pP_p = \tfrac12\sum_c(F_n\lambda + F_t\eta)$, que no
    depende del área.
\end{itemize}

\end{document}
